\documentclass[12pt]{article}

% Layout.
\usepackage[top=1in, bottom=0.75in, left=1in, right=1in, headheight=1in, headsep=6pt]{geometry}

% Fonts.
\usepackage{mathptmx}
\usepackage[scaled=0.86]{helvet}
\renewcommand{\emph}[1]{\textsf{\textbf{#1}}}

% TiKZ.
\usepackage{tikz, pgfplots}
\usetikzlibrary{calc}

% Misc packages.
\usepackage{amsmath,amssymb,latexsym}
\usepackage{graphicx}
\usepackage{array}
\usepackage{xcolor}
\usepackage{multicol}
\usepackage{fancyhdr}
\pagestyle{fancy}

\usepackage[parfill]{parskip}


% Misc.
\renewcommand{\d}{\displaystyle}
\newcommand{\ds}{\displaystyle}
\newcommand{\ul}[1]{\underline{#1}}
\def\bc{\begin{center}}
\def\ec{\end{center}}
\def\be{\begin{enumerate}}
\def\ee{\end{enumerate}}

\newcommand{\ans}[1][1in]{\rule{#1}{.5pt}}


\lhead{Math F113X: Exam I Review}
\rhead{}

\begin{document}
%\begin{center} {\Large{Voting Theory}} \\ \quad \\You should know \end{center}

{\color{red} Note that before Fall 2026, Exam 1 covered Voting Theory, Weighted Voting, and Fair Division. This Exam 1 review material covers Voting Theory and Weighted Voting only.}

\section{Voting Theory}

\subsection{You should know} 

\noindent\textbf{Terminology:} majority, plurality, Cordorcet Winner

\noindent{\textbf{Voting Methods by Name}}: Plurality Method, Instant Run-off Voting (IRV) (also known as Ranked Choice Voting), Borda Count, Copeland's Method

\noindent{\textbf{Fairness Criteria by Name}}: Condorcet Criterion, Monotonicity Criterion, Majority Criterion, Irrelevance of Independent Alternatives (IIA) Criterion

%\begin{center} {\Large{Sample Voting Theory Problems}}  \end{center}
\subsection{Sample Voting Theory Problems}

\be

\item A class of middle schoolers are trying to decide what would be the worst ingredient to add to a cake their Principal has to eat if they win the Science Bowl. They narrow their options down to three: pickled onions (PO), stinky cheese (SC), or anchovies (A). The class is polled and the resulting preference schedule is below.	\\

\begin{center}
\begin{tabular}{l || c|c|c|c|c|c}
&22&5&23&2&11&20\\
\hline \hline
1st choice& PO&PO&SC&SC&A&A\\
\hline
2nd choice&SC&A&PO&A&PO&SC\\
\hline
3rd choice&A&SC&A&PO&SC&PO
\end{tabular}
\end{center}
\begin{enumerate}
	\item How many students voted? \ans
	
	\bigskip
	\item How many votes does a candidate need to win in order to win a \text{majority}? Show the calculation that gives your answer.
	\ans
	
	\bigskip

	\item How many votes does a candidate need to win in order to win a \text{plurality}? Show the calculation that gives your answer.
	\ans
	
	\bigskip

	\item Find the winner using the plurality method. Show your work.
	\vfill
	\item Find the winner using Instant Runoff Voting.  Show your work.
	\vfill
	
	\vfill
	\newpage
	
	For reference, here's the preference schedule again. 
	
	\begin{center}
\begin{tabular}{l || c|c|c|c|c|c}
&22&5&23&2&11&20\\
\hline \hline
1st choice& PO&PO&SC&SC&A&A\\
\hline
2nd choice&SC&A&PO&A&PO&SC\\
\hline
3rd choice&A&SC&A&PO&SC&PO
\end{tabular}
\end{center}
	\item Find the winner using Borda Count.  Show your work.
	\vfill
	\item Find the winner using Copeland's Method. Show your work.
	\vfill
	\item Is any candidate a Condorcet Winner? Explain your answer.	
	\vfill
	\end{enumerate}
	
	\newpage
	
	\item Consider the following preference schedule, for an election with five candidates: Alfred, Bethanne, Carol, David, and Edgar.
	
	\begin{tabular}{c|cccccc}
\textbf{Rank} & \textbf{40} & \textbf{8} & \textbf{25} & \textbf{15} & \textbf{5} & \textbf{7} \\
\hline
1st & A & B & C & D & E & C \\
2nd & C & A & A & A & A & B \\
3rd & B & C & B & B & B & A \\
4th & D & D & D & C & C & D \\
5th & E & E & E & E & D & E \\
\end{tabular}

 Find the winner under Instant Runoff Voting. Clearly show your rounds, and state who is eliminated in each round. For each round, clearly indicate the vote total for each candidate.
	
	

\vfill

The winner is \ans and they won in \ans rounds.

\bigskip

\item Consider the following preference schedule, and the points assigned under Borda count. Explain which fairness criterion this shows Borda Count fails, and why this particular election fails that fairness criterion.

\begin{tabular}{c||c|c|c}
\textbf{Rank} &  \textbf{51} & \textbf{30} & \textbf{19} \\
\hline
1st &  A & B & C \\
2nd &  B & C & D \\
3rd & C & D & B \\
4th & D & A & A \\
\end{tabular}
\hspace{1cm}
\begin{tabular}{c | c}
Candidate & Borda Count Points \\ \hline
A & 253\\
B & 311 \\
C & 268\\
D & 168
\end{tabular}

	\vspace{1in}
	
	\newpage
	
\item For each of the following fairness criteria, describe what an election would look like that \emph{failed} that criterion. (You do not need to provide an actual preference schedule, just talk about what the problem would be.)	
\be
\item Condorcet Criterion

\vfill

\item Monotonicity Criterion

\vfill

\item Majority Criterion

\vfill

\item Irrelevance of Independent Alternatives (IIA) Criterion

\vfill

\ee
\ee

\newpage

%\begin{center} {\Large{Weighted Voting}} \\ \end{center}
\section{Weighted Voting}

\subsection{Terminology} quota, weight, winning coalition, critical player in a winning coalition, dictator, veto power, dummy, Banzhaf Power Index or Banzhaf Power Distribution.

%\begin{center} {\Large{Sample Weighted Voting Problems}}  \end{center}

\subsection{Sample Weighted Voting Problems}
\begin{enumerate}
\item Nova Robotics Inc. has five board members who vote on major company decisions (funding rounds, executive hires, acquisitions, etc.). Each member's voting weight is based on the percentage of company shares they control, and company bylaws require any motion to pass with at least 60\% of the total votes.

Founder/CEO: 56 shares\\
Co-founder/CTO: 44 shares\\
Lead Venture Capital firm: 40 shares\\
Early Angel Investor: 36 shares\\
Employee Stock Pool Representative: 24 shares\\

\be
\item Write a weighted voting system of the form $[q: w_{1}, \ldots, w_{n}]$ that represents this scenario. 

\vspace{1in}

\item Find a winning coalition in this system, and write a sentence explaining how you know it is a winning coalition.


\vspace{1in}
\ee

\newpage
\item For each weighted voting system below, determine if there are any dictators, anyone with veto power, or any dummies.
	\begin{enumerate}
	\item $[16:16,11,3,1]$
	\vfill
	\item $[51:40,30,20,10]$
	\vfill
	\item $[31:10,9,8,7,6]$
	\vfill
	\end{enumerate}
	
	\newpage
\item Consider the weighted voting system $[q:21,20,7,5]$ where the players can pass a motion with a \emph{majority}.
	\begin{enumerate}
	\item What is $q$ in this case? How do you know?
	
	\bigskip
		\bigskip

	
	\item List all winning coalitions.
	\vfill
	\item In each of the winning coalitions above, indicate who is a critical player.
	
	\item Calculate the Banzhaf Power Index for each player.
	\vfill
	\item Based on your calculations above, are there any dummies? Explain.
	\vfill
	\end{enumerate}
\end{enumerate}
\newpage	
%%\begin{center} {\Large{Fair Division}}  \end{center}
%\section{Fair Division}
%
%%\noindent\textbf{Methods by Name:} 
%\subsection{Methods by Name}
%Divider-Chooser, Lone Divider, the Method of Sealed Bids.
%
%%\begin{center} {\Large{Sample Fair Divison Problems}}  \end{center}
%\subsection{Sample Fair Divison Problems}
%
%\begin{enumerate}
%\item For each of the three methods listed above, what is the \textbf{context} in which they are appropriate? (for example: how many people are dividing? How can you split up the stuff?)
%
%\vfill
%
%\item The Cookie Monster and Elmo are splitting a giant cookie that is half chocolate-chip and half oatmeal, worth \$12. The Cookie Monster likes chocolate chip twice as much as oatmeal. Elmo likes oatmeal three times as much as chocolate chip. 
%
%\be
%\item Suppose someone splits the cookie with all chocolate chip on one side and all oatmeal on the other. Attach dollar values to each share for each person.
%
%
%\begin{tabular}{l | c | c}
%&share 1: chocolate chip &share 2: oatmeal \\
%\hline
%&&\\
%Cookie Monster&&\\
%&&\\
%\hline
%&&\\
%Elmo&&\\ &&\\
%\end{tabular}
%
%\vfill
%
%\item Which cookie pieces (if any) are fair shares for Elmo? Which are fair shares for Cookie Monster? Show calculations that support your answer.
%
%\vfill
%\ee
%\newpage
%
%\item
%\begin{enumerate}
%\item James, Gordon, Julie, Alexei, and Latrice divide their pile of Halloween candy worth a total of \$20. Determine who was the divider and determine how the lone divider method would proceed.\\
%
%\begin{tabular}{l || c | c | c| c| c}
%&Pile 1&Pile 2&Pile 3&Pile 4&Pile 5\\
%\hline \hline
%James&\$0&\$5&\$10&\$5&\$0\\ \hline
%Gordon&\$2&\$4&\$6&\$2&\$6\\ \hline
%Julie&\$4&\$4&\$4&\$4&\$4\\ \hline
%Alexei&\$2&\$4&\$12&\$2&\$0\\ \hline
%Latrice&\$2&\$10&\$7&\$1&\$0\\ 
%\end{tabular}
%\vfill
%	\item Suppose \textbf{Gordon} changes his evaluation of the piles. Now determine how the lone divider method would proceed.\\
%	\begin{tabular}{l || c | c | c| c| c}
%&Pile 1&Pile 2&Pile 3&Pile 4&Pile 5\\
%\hline \hline
%James&\$0&\$5&\$10&\$5&\$0\\ \hline
%\textbf{Gordon}&\$2&\textbf{\$7}&\$6&\$2&\textbf{\$3}\\ \hline
%Julie&\$4&\$4&\$4&\$4&\$4\\ \hline
%Alexei&\$2&\$4&\$12&\$2&\$0\\ \hline
%Latrice&\$2&\$10&\$7&\$1&\$0\\ 
%\end{tabular}
%	\end{enumerate}
%\vfill
%\newpage
%
%
%	
%	
%\item Harry Potter (P), Hermione Granger (G), Luna Lovegood (L), and Ronald Weasley (W) are dividing up some loot they found in the dungeons. The loot consists of a velvet cloak, a gold chalice, and a self-cleaning cauldron. They decide to divide the loot using the sealed bid method. The table below shows how many gold coins each person bid for each item.\\
%
%\begin{tabular}{l || c|c|c|c}
%&cloak&chalice&cauldron\\
%\hline \hline
%Potter&60 gold coins&60 gold coins&0 gold coins\\
%\hline
%Granger&50 gold coins&10 gold coins&80 gold coins\\
%\hline
%Lovegood&100 gold coins&0 gold coins&100 gold coins\\
%\hline
%Weasley&30 gold coins&20 gold coins& 30 gold coins\\
%\end{tabular}
%	\begin{enumerate}
%	\item Determine each person's fair share.
%	\vfill
%	\item Determine which person gets each item.
%	\vfill
%	\item Determine how many gold coins each of them owes to the holding pile or receives from the holding pile.
%	
%	\vfill
%%	\newpage
%	\item Determine the surplus. (Fractional gold coins are fine.)
%	\vfill
%	\item Determine the final allotment.
%	\vfill
%	\end{enumerate}
%\end{enumerate}
\end{document}